\section{Retained numerical producer and new inspection outputs}\label{sec:r9-executor}
\subsection{Immutable source, checkpoint policy and causal replay}
The retained airframe-relative state is $x=(t,R,b,m_t,m_R)$ with an 18-dimensional local covariance. Measured motion increments propagate pose, bias and covariance. A scan and its source-time prior form one immutable source. Robust point-to-plane registration \cite{chen1992,pomerleau2015} exposes checkpoints at up to 1, 2, 4, 8, 16 and 25 iterations. Each candidate uses the same scan/prior; later versions replace earlier versions rather than introduce another independent likelihood factor. At an already acquired receiver time $t_d$, a delivered candidate is assigned after replay of the measured increments:
\begin{equation}
(x_{t_d}^{(j)},P_{t_d}^{(j)})=\mathcal R_{s\rightarrow t_d}(x_s^{(j)},P_s^{(j)};u_{s+1:t_d}).\label{eq:r9-replay}
\end{equation}
Ordered rollback installs the source prior after a rejected completion. Original source/prior seals, producer records, version ordering and history checks are retained. Their source-replacement and exposure proofs, detailed solver objective and original experiments remain in the supplementary and archival material.

The frozen CART and random-forest forecasts predict future endpoint error, event frequency and remaining work. Numerical selection trades predicted four-second error, failure cost and computation, with common stage-budget reservation. The primary new catalogue includes Fixed2, Fixed4, Final, FlatCART, FlatRF, KG1, ContractRF and OldRAVC. All receive the same current source, measured increments, numerical predictors, prices, budget wrapper and service load. KG1 adds the retained time/profile predicates; its name does not denote a learned graph embedding. The predictors are not refitted to the new camera tasks. The stronger procedural evidence baseline receives all task knowledge, rather than an artificially weakened numerical producer.

One actual common solver prefix is executed for each synthetic input. The separately executed policy/service wrappers consume only their revealed checkpoints and charge the corresponding prefix. This controlled trace design is paired counterfactual execution, not 32 independent inputs or 32 separately timed complete cold solves. Native component time and the deterministic work clock have different measurement boundaries. Historical full-route online timings remain a separate evidence block.

\subsection{Actual raster and diagnostic target processing}
A prescribed nominal route and prior CAD mesh plan up to 15 face/barycentric target addresses before stochastic map displacement and sensing. A synthetic gimbal aims at a prior-CAD patch. Targets remain the same across the paired ROI intervention. Raycasting against the perturbed geometric scene creates RGB pixels with procedural texture, fixed lighting, projected 25-mm-radius diagnostic landmarks, an exposure-blur surrogate and Gaussian intensity noise. A nearest-palette decoder estimates landmark pixel centers using only the image and known landmark colors. Ground truth is used for scene generation and offline correctness labels, not for decoding or online numerical policy choice. This is a diagnostic synthetic landmark task, not crack detection or a new real-camera experiment.

At a target with range $d$ and absolute ray/normal incidence $\gamma$, the defined geometric sampling density is $g=d/(f_{\rm px}\gamma)$. Fixed-target quality uses the maximum GSD and motion/exposure blur across the preplanned addresses, including targets outside the current ROI. Correctly acquired target coverage is scored from actual decoded pixels. The capture-time geotag projects the CAD targets using the causally installed localization state. Its pixel discrepancy and metric landmark displacement are separately scored. An absent positive source version fails the readiness requirement; a later update cannot retroactively count as capture-time ready.

Image capture occurs at an independent, before-sensing sampled tick $c\in\{0.1,0.2,0.3\}$~s. Report delay and image age are distinct:
\begin{equation}
t_{\rm report}=\max\{t_{\rm first\ positive},c+0.05+\alpha w_{\rm camera}\},\qquad
a_{\rm image}=t_{\rm report}-c.\label{eq:r9-report}
\end{equation}
The camera clock characterizes the executed \emph{virtual raster producer plus decoder}, not a physical camera's capture latency. Its excluded-pilot median is frozen separately for each ROI. The same declared load multiplier is used across methods. Native raster generation, decoding and localization components are reported separately; their sum is a mixed-boundary virtual workload, not embedded online CPU or worst-case timing. This narrow timing model is part of the certified simulation event and cannot be transferred directly to physical inspection.
